\subsection{METHODS OF PROOF}

I. Method of the auxiliary hypothesis. This rests on the following rule:

C14 (Criterion of deduction). Let A be a relation in C, and let \(C'\) be the theory obtained by adjoining A to the axioms of C. If B is a theorem in \(C'\) , then \(A \Longrightarrow B\) is a theorem in C.

Let \(B_1, B_2, \ldots, B_n\) be a proof in \(\mathfrak{C}'\) in which \(B\) appears. We shall show, step by step, that the relations \(A \Longrightarrow B_k\) are theorems in \(\mathfrak{C}\) . Suppose that this has been established for the relations which precede \(B_i\) , and let us show that \(A \Longrightarrow B_i\) is a theorem in \(\mathfrak{C}\) . If \(B_i\) is an axiom of \(\mathfrak{C}'\) , then \(B_i\) is either an axiom of \(\mathfrak{C}\) or is \(A\) . In both cases, \(A \Longrightarrow B_i\) is a theorem in \(\mathfrak{C}\) by applying C9 or C8. If \(B_i\) is preceded by relations \(B_j\) and \(B_j \Longrightarrow B_i\) , we know that \(A \Longrightarrow B_j\) and \(A \Longrightarrow (B_j \Longrightarrow B_i)\) are theorems in \(\mathfrak{C}\) . Hence \((B_j \Longrightarrow B_i) \Longrightarrow (A \Longrightarrow B_i)\) is a theorem in \(\mathfrak{C}\) by C13. Hence, by C6, \(A \Longrightarrow (A \Longrightarrow B_i)\) , that is to say ``(not \(A\) ) or \((A \Longrightarrow B_i)\) '', is a theorem in \(\mathfrak{C}\) , and therefore so is `` \((A \Longrightarrow B_i)\) or (not \(A\) )'' by S3. Now, \((\text{not } A) \Longrightarrow ((\text{not } A) \text{ or } B_i)\) , that is to say \((\text{not } A) \Longrightarrow (A \Longrightarrow B_i)\) , is a theorem in \(\mathfrak{C}\) , by S2. By application of S4 we see then that

\[
((\mathbf {A} \Rightarrow \mathbf {B} _ {i}) \text {   or   (not   } \mathbf {A})) \Rightarrow ((\mathbf {A} \Rightarrow \mathbf {B} _ {i}) \text {   or   } (\mathbf {A} \Rightarrow \mathbf {B} _ {i}))
\]

is a theorem in \(\mathfrak{C}\) , and hence that `` \((A \Longrightarrow B_i)\) or \((A \Longrightarrow B_i)\) '' is a theorem in \(\mathfrak{C}\) . By S1 we conclude that \(A \Longrightarrow B_i\) is a theorem in \(\mathfrak{C}\) .

In practice we indicate that we are going to use this criterion by a phrase such as ``suppose that A is true''. This phrase means that for the time being the reasoning will be performed in the theory \(C'\) , until the relation B has been proved. When this has been achieved it has been established that \(A \Longrightarrow B\) is a theorem in C, and one continues thereafter to reason in C without in general indicating that one has abandoned the theory \(C'\) . The relation A introduced as a new axiom is called the auxiliary hypothesis. * For example, when we say ``let x be a real number'', we are constructing a theory in which the relation ``x is a real number'' is an auxiliary hypothesis. *

\begin{enumerate}
\def\labelenumi{\Roman{enumi}.}
\setcounter{enumi}{1}
\tightlist
\item
  Method of reductio ad absurdum. This is founded on the following rule :
\end{enumerate}

C15. Let A be a relation in C, and let \(C'\) be the theory obtained by adjoining the axiom ``not A'' to the axioms of C. If \(C'\) is contradictory, then A is a theorem in C.

For A is a theorem in \(C'\) ; consequently (method of the auxiliary hypothesis) ``(not \(A\) ) \(\Longrightarrow\) \(A\) '' is a theorem in C. By S4,

\[
(A \text {   or   (not   } A)) \Longrightarrow (A \text {   or   } A)
\]

is a theorem in \(\mathfrak{C}\) ; by C10, `` \(A\) or \(A\) '' is a theorem in \(\mathfrak{C}\) . Now use S1.

In practice, we indicate that we are going to use this criterion by a phrase such as ``suppose that A is false''. This phrase means that for the time being the reasoning will be performed in the theory \(C'\) , until two theorems of the form B and ``not B'' have been proved. When this has been achieved it is then established that A is a theorem in C, which is generally indicated by a phrase such as ``Now this (i.e., in the preceding notation, B and ``not B'') is absurd; hence A is true''. One then continues in the original theory C.

¶ As first applications of these methods, let us establish the following criteria:

C16. If \(A\) is a relation in \(\mathfrak{C}\) , then (not not \(A\) ) \(\Longrightarrow A\) is a theorem in \(\mathfrak{C}\) .

For suppose that ``not not A'' is true; then we have to prove A. Suppose A is false. In the theory so defined, ``not not A'' and ``not A'' are theorems, which is absurd; therefore A is true.

C17. If A and B are relations in \(\mathfrak{G}\) , then

\[
((\text { not } B) \Longrightarrow (\text { not } A)) \Longrightarrow (A \Longrightarrow B)
\]

is a theorem in \(\mathfrak{C}\) .

For suppose that \((\text{not } B) \Longrightarrow (\text{not } A)\) is true. We have to show that \(A \Longrightarrow B\) is true. Suppose that A is true, and let us show that B is true. Suppose ``not \(B\) '' is true. Then ``not A'' is true, which is absurd.

\begin{enumerate}
\def\labelenumi{\Roman{enumi}.}
\setcounter{enumi}{2}
\tightlist
\item
  Method of disjunction of cases. This rests on the following rule:
\end{enumerate}

C18. Let \(A, B, C\) be relations in \(\mathfrak{G}\) . If `` \(A\) or \(B\) '', \(A \Longrightarrow C\) , and \(B \Longrightarrow C\) are theorems in \(\mathfrak{G}\) , then \(C\) is a theorem in \(\mathfrak{G}\) .

For, by S4, ``(A or B)⇒ (A or C)'' and ``(C or A)⇒ (C or C)'' are theorems in C. By S3 and S1, it follows that (A or B)⇒ C is a theorem in C; hence the result.

To prove C it is therefore enough, when we have at our disposal a theorem ``A or B'', first to prove C by adjoining A to the axioms of G, and then to prove C by adjoining B to the axioms of G. The interesting feature of this method lies in the fact that if ``A or B'' is true, we cannot in general assert either that A is true or that B is true.

In particular, by C10, if `` \(A \Rightarrow C\) '' and ``(not \(A) \Rightarrow C\) '' are both theorems in \(\mathfrak{C}\) , then \(C\) is a theorem in \(\mathfrak{C}\) .

\begin{enumerate}
\def\labelenumi{\Roman{enumi}.}
\setcounter{enumi}{3}
\tightlist
\item
  Method of the auxiliary constant. This is founded on the following rule:
\end{enumerate}

C19. Let x be a letter and let A and B be relations in \(\mathfrak{G}\) such that :

\begin{enumerate}
\def\labelenumi{(\arabic{enumi})}
\item
  the letter \(\pmb{x}\) is not a constant of \(\mathfrak{C}\) and does not appear in \(\pmb{B}\) ;
\item
  there is a term \(\pmb{T}\) in \(\mathfrak{C}\) such that \((\pmb{T}|\pmb{x})\pmb{A}\) is a theorem in \(\mathfrak{C}\) .
\end{enumerate}

Let \(\mathfrak{C}'\) be the theory obtained by adjoining \(A\) to the axioms of \(\mathfrak{C}\) . If \(B\) is a theorem in \(\mathfrak{C}'\) , then \(B\) is a theorem in \(\mathfrak{C}\) .

Indeed, \(A \Longrightarrow B\) is a theorem in \(\mathcal{C}\) (criterion of deduction). Since \(x\) is not a constant of \(\mathcal{C}\) , \((T|x)(A \Longrightarrow B)\) is a theorem in \(\mathcal{C}\) by virtue of C3. Since \(x\) does not appear in \(B\) , \((T|x)(A \Longrightarrow B)\) is identical with \(((T|x)A) \Longrightarrow B\) , by CS5 (§ 1, no. 2). Finally, \((T|x)A\) is a theorem in \(\mathcal{C}\) , and therefore so is \(B\) .

Intuitively, the method consists in using, in order to prove B, an arbitrary object x (the auxiliary constant) which is supposed to be endowed with certain properties, denoted by A. * For example, in a proof in geometry which involves, among other things, a line D, we may ``take'' a point x on this line; the relation A is then \(x \in D\) . * In order that one should be able to use an object endowed with certain properties during the course of a proof, it is clearly necessary that such objects should exist. The theorem \((T|x)A\) , called the theorem of legitimation, guarantees this existence.

In practice we indicate that we are going to use this method by a phrase such as ``let \(\pmb{x}\) be an object such that \(A\)''. By contrast with the method of the auxiliary hypothesis, the conclusion of the argument does not involve \(\pmb{x}\) .

